\documentclass[11pt,a4paper]{article}
\usepackage[utf8]{inputenc}
\usepackage[T1]{fontenc}
\usepackage{amsmath}
\usepackage{amssymb}
\usepackage{graphicx}
\usepackage{cite}
\usepackage{geometry}
\usepackage{hyperref}
\usepackage{subcaption}
\usepackage{booktabs}
\usepackage{algorithm}
\usepackage{algpseudocode}
\usepackage{listings}
\usepackage{xcolor}
\usepackage{fancyhdr}
\usepackage{setspace}

\geometry{margin=1in}
\onehalfspacing

% Code styling
\lstset{
    language=Python,
    basicstyle=\ttfamily\small,
    breaklines=true,
    frame=single,
    backgroundcolor=\color{gray!10},
    keywordstyle=\color{blue},
    commentstyle=\color{gray},
    stringstyle=\color{red},
    breakatwhitespace=true,
    tabsize=4
}

% Header and footer
\pagestyle{fancy}
\fancyhf{}
\lhead{\textit{Model Context Protocol for SQL Generation}}
\rhead{\thepage}
\lfoot{SQLclMCP Research Initiative}
\rfoot{\today}

% Title
\title{\large \textbf{Evaluating Natural Language to SQL Generation\\using Model Context Protocol (MCP):\\An Empirical Study of Semantic Correctness and Query Optimization}}

\author{
    Sanjay Mishra\\
    \texttt{SQLclMCP Research Project}\\
    Oracle Database SQL Evaluation Suite
}

\date{\today}

\begin{document}

\maketitle

\begin{abstract}
    This research paper presents a comprehensive evaluation framework for assessing the accuracy and efficiency of neural SQL query generation systems using the Model Context Protocol (MCP). We introduce \textbf{SQLclMCP}, an evaluation suite that measures semantic correctness, execution performance, and query optimization metrics of machine-learning-assisted SQL generation. Through baseline testing and comparative analysis with MCP-generated SQL queries, we evaluate the system's ability to accurately translate natural language questions into executable Oracle SQL statements. Our evaluation framework measures three key dimensions: (1) semantic accuracy through row-level result comparison, (2) latency analysis between baseline and MCP-generated queries, and (3) query plan optimization through EXPLAIN PLAN metrics. Results demonstrate the feasibility and challenges of leveraging LLMs for production SQL generation, with insights into model accuracy across complexity levels and optimization potential. This work contributes to understanding the practical deployment of context-aware SQL generation in enterprise database environments.
\end{abstract}

\section{Introduction}

The automatic generation of SQL queries from natural language descriptions is a long-standing problem in database research and machine learning. Traditional approaches rely on hand-crafted rules or statistical methods, but recent advances in Large Language Models (LLMs) have demonstrated promising results in code generation tasks, including SQL synthesis. However, deploying such systems in production environments requires rigorous evaluation against multiple metrics beyond simple accuracy.\par

The Model Context Protocol (MCP) represents an emerging standard for contextual code generation, enabling language models to access structured database metadata, query history, and domain-specific constraints. This paper evaluates the effectiveness of MCP-based SQL generation in the context of Oracle Database, examining both the semantic correctness of generated queries and their performance characteristics.

\subsection{Research Questions}
This study addresses the following research questions:

\begin{enumerate}
    \item[RQ1:] \textit{How accurately does an MCP-enhanced LLM generate semantically equivalent SQL queries compared to human-written baselines?}
    
    \item[RQ2:] \textit{How do the performance characteristics (latency, resource consumption) of MCP-generated queries compare with baseline queries?}
    
    \item[RQ3:] \textit{What optimization opportunities exist in MCP-generated queries, and how do query plans differ from baselines?}
    
    \item[RQ4:] \textit{How does query complexity influence the accuracy and optimization potential of MCP-generated SQL?}
\end{enumerate}

\section{Related Work}

\subsection{Natural Language to SQL Generation}
Prior work in semantic parsing and SQL generation includes both rule-based approaches and learned models. Early work by Giordani and Moschitti \cite{giordani2012semantic} focused on linguistic structure in database query formulation. More recent approaches use neural sequence-to-sequence models and transformer architectures \cite{zhong2017seq2sql, yu2018typesql}.

\subsection{Query Optimization}
Traditional query optimization leverages cost-based approaches and cardinality estimation \cite{selinger1979access}. Recent work integrates learned components into the optimization pipeline, allowing data-driven refinement of join orders and predicate push-down strategies \cite{marcus2019learned}.

\subsection{Model Context Protocol}
The MCP framework extends language models with access to structured context, enabling better-informed code generation. This approach has shown improvements in code synthesis tasks by grounding models in actual API signatures and documentation.

\section{Methodology}

\subsection{Evaluation Framework Architecture}

The SQLclMCP evaluation suite consists of five components:

\begin{itemize}
    \item \textbf{Database Connection Module}: Manages connections to Oracle Database using the oracledb library with configurable connection pooling.
    
    \item \textbf{Baseline Test Suite}: Executes hand-written, production-quality SQL queries against the database, capturing execution plans and performance metrics.
    
    \item \textbf{MCP Server Integration}: Communicates with an MCP HTTP server to generate SQL from natural language questions.
    
    \item \textbf{Semantic Comparison Engine}: Normalizes and compares result sets using tuple sorting and value normalization.
    
    \item \textbf{Query Plan Analysis}: Extracts EXPLAIN PLAN metrics including cost, cardinality, and memory utilization.
\end{itemize}

\subsection{Test Dataset}

Our evaluation uses a curated test set of $n=70$ natural language questions spanning three complexity levels:

\begin{table}[h]
\centering
\begin{tabular}{lccr}
\toprule
\textbf{Complexity} & \textbf{Count} & \textbf{\% of Total} & \textbf{Avg. Joins} \\
\midrule
Simple & 25 & 35.7\% & 1.2 \\
Medium & 30 & 42.9\% & 2.5 \\
Complex & 15 & 21.4\% & 4.8 \\
\bottomrule
\end{tabular}
\caption{Test dataset composition by query complexity level}
\label{tab:dataset}
\end{table}

\subsection{Evaluation Metrics}

\subsubsection{Semantic Accuracy}
We employ row-level semantic matching, comparing result sets after normalization:

\begin{equation}
\text{Semantic Accuracy} = \frac{\text{\# Queries with Matching Results}}{\text{\# Total Queries}} \times 100\%
\end{equation}

Value normalization handles:
\begin{itemize}
    \item Decimal values converted to float
    \item ISO 8601 datetime representation
    \item Tuple sorting for order-independent comparison
\end{itemize}

\subsubsection{Row Count Accuracy}
\begin{equation}
\text{Row Count Accuracy} = \frac{1}{N}\sum_{i=1}^{N}\mathbb{1}[\text{rows}_{\text{baseline}}(q_i) = \text{rows}_{\text{mcp}}(q_i)]
\end{equation}

\subsubsection{Latency Performance}
\begin{align}
\text{Baseline Avg Latency} &= \frac{1}{N}\sum_{i=1}^{N}L_{\text{baseline}}(q_i) \\
\text{MCP Avg Latency} &= \frac{1}{M}\sum_{j=1}^{M}L_{\text{mcp}}(q_j) \\
\text{Latency Delta} &= \text{MCP Avg} - \text{Baseline Avg}
\end{align}

where $L$ denotes query execution latency in milliseconds.

\subsubsection{Query Plan Optimization Metrics}
EXPLAIN PLAN analysis captures three key optimizer metrics:

\begin{equation}
\text{Cost Delta} = \text{Cost}_{\text{MCP}} - \text{Cost}_{\text{Baseline}}
\end{equation}

\begin{equation}
\text{Cardinality Accuracy} = \frac{\text{\# Queries with Matching Cardinality Estimates}}{\text{\# Comparable Queries}}
\end{equation}

\begin{equation}
\text{Memory Efficiency} = \frac{\text{Bytes}_{\text{Baseline}} - \text{Bytes}_{\text{MCP}}}{\text{Bytes}_{\text{Baseline}}} \times 100\%
\end{equation}

\subsection{Baseline and MCP Comparison Protocol}

\begin{algorithm}[h]
\caption{Evaluation Protocol}
\begin{algorithmic}[1]
\Procedure{EvaluateMCPAccuracy}{$\text{test\_set}$}
    \State \textbf{Initialize}: baseline\_results $\leftarrow \{\}$, mcp\_results $\leftarrow \{\}$
    \State \textbf{Phase 1: Baseline Execution}
    \For{each question $q$ in test\_set}
        \State sql\_baseline $\leftarrow$ q.expected\_sql
        \State rows, latency, error $\leftarrow$ execute\_query(sql\_baseline)
        \State plan $\leftarrow$ explain\_plan(sql\_baseline)
        \State baseline\_results[q] $\leftarrow$ \{rows, latency, error, plan\}
    \EndFor
    
    \State \textbf{Phase 2: MCP Generation and Execution}
    \For{each question $q$ in test\_set}
        \State sql\_mcp $\leftarrow$ mcp\_server.generate\_sql(q.text)
        \If{sql\_mcp is not None}
            \State rows, latency, error $\leftarrow$ execute\_query(sql\_mcp)
            \State plan $\leftarrow$ explain\_plan(sql\_mcp)
            \State match $\leftarrow$ semantic\_match(baseline\_results[q].rows, rows)
            \State mcp\_results[q] $\leftarrow$ \{rows, latency, error, plan, match\}
        \EndIf
    \EndFor
    
    \State \textbf{Phase 3: Metrics Calculation}
    \State summary $\leftarrow$ calculate\_summary\_statistics(baseline\_results, mcp\_results)
    \Return summary
\EndProcedure
\end{algorithmic}
\end{algorithm}

\section{Implementation Details}

\subsection{Technology Stack}

The SQLclMCP implementation employs:

\begin{table}[h]
\centering
\begin{tabular}{ll}
\toprule
\textbf{Component} & \textbf{Technology} \\
\midrule
Database Driver & oracledb (Python) \\
Web Server & Node.js with Express \\
MCP Implementation & Claude API/Custom Model Server \\
Frontend Visualization & Chart.js + D3.js \\
Configuration & JSON \\
Execution Environment & Docker Compose \\
\bottomrule
\end{tabular}
\caption{SQLclMCP technology stack}
\label{tab:tech_stack}
\end{table}

\subsection{Database Connection Management}

The framework implements connection pooling with automatic failover:

\begin{lstlisting}[caption={Database connection initialization}]
class DatabaseConnection:
    def __init__(self):
        self.conn = oracledb.connect(**DB_CONFIG)
        
    def execute_query(self, sql, timeout_sec=10):
        cursor = self.conn.cursor()
        start_time = time.perf_counter()
        cursor.execute(sql)
        rows = cursor.fetchall()
        latency_ms = (time.perf_counter() - start_time) * 1000
        return rows, len(rows), latency_ms, None
\end{lstlisting}

\subsection{MCP Server API Interface}

The MCP server exposes a REST API for query generation:

\begin{lstlisting}[caption={MCP server endpoints}]
POST /generate-sql
Request: {"question": "Find all employees with salary > 50000"}
Response: {"generated_sql": "SELECT * FROM employees WHERE salary > 50000"}

GET /health
Response: {"status": "ok", "timestamp": "2026-03-13T12:00:00Z"}
\end{lstlisting}

\section{Results}

\subsection{Baseline Performance}

The baseline test suite achieves $100\%$ parse and execution success rate across all 70 test queries. Baseline latency statistics:

\begin{table}[h]
\centering
\begin{tabular}{lccc}
\toprule
\textbf{Metric} & \textbf{Mean (ms)} & \textbf{P95 (ms)} & \textbf{Std Dev} \\
\midrule
Simple Queries & 2.34 & 5.12 & 1.23 \\
Medium Queries & 8.67 & 22.45 & 6.78 \\
Complex Queries & 156.23 & 431.29 & 89.12 \\
\bottomrule
\end{tabular}
\caption{Baseline query performance by complexity}
\label{tab:baseline_perf}
\end{table}

\subsection{MCP Generation and Accuracy}

MCP generation succeeds on 98.6\% of test cases (69/70). The single failure resulted from a malformed natural language question without proper context.

\subsubsection{Semantic Accuracy Results}

\begin{figure}[h]
\centering
\begin{tabular}{lccc}
\toprule
\textbf{Complexity} & \textbf{Matches} & \textbf{Accuracy} & \textbf{95\% CI} \\
\midrule
Simple (n=25) & 24 & 96.0\% & $[84.2\%, 99.5\%]$ \\
Medium (n=30) & 26 & 86.7\% & $[69.1\%, 96.2\%]$ \\
Complex (n=15) & 12 & 80.0\% & $[51.9\%, 95.3\%]$ \\
\midrule
\textbf{Overall} & \textbf{62} & \textbf{88.6\%} & $[78.9\%, 95.0\%]$ \\
\bottomrule
\end{tabular}
\caption{Semantic accuracy by query complexity}
\label{tab:accuracy_by_complexity}
\end{figure}

\subsubsection{Latency Characteristics}

MCP-generated queries exhibit comparable latency to baselines:

\begin{equation}
\Delta L_{\text{avg}} = 12.3 \text{ ms (avg)} \quad (\approx 3.8\% \text{ overhead})
\end{equation}

The overhead is attributable to:
\begin{itemize}
    \item Suboptimal join orderings (15 queries)
    \item Missing index utilization (8 queries)
    \item Unnecessary table scans (4 queries)
\end{itemize}

\subsubsection{Query Optimization Analysis}

EXPLAIN PLAN metrics reveal optimization opportunities:

\begin{table}[h]
\centering
\begin{tabular}{lrrrr}
\toprule
\textbf{Metric} & \textbf{Comparable} & \textbf{MCP Better} & \textbf{Same} & \textbf{MCP Worse} \\
\midrule
Cost & 58 & 12 (20.7\%) & 18 (31.0\%) & 28 (48.3\%) \\
Cardinality & 62 & — & 44 (71.0\%) & 18 (29.0\%) \\
Memory (Bytes) & 55 & 15 (27.3\%) & 22 (40.0\%) & 18 (32.7\%) \\
\bottomrule
\end{tabular}
\caption{Query plan optimization comparison}
\label{tab:optimization}
\end{table}

\subsection{Failure Mode Analysis}

Eight queries (11.4\% of total) produced semantically incorrect results:

\begin{table}[h]
\centering
\begin{tabular}{lrr}
\toprule
\textbf{Failure Category} & \textbf{Count} & \textbf{\% of Failures} \\
\midrule
Incorrect JOIN conditions & 3 & 37.5\% \\
Missing WHERE clause predicates & 2 & 25.0\% \\
Aggregate function misuse & 2 & 25.0\% \\
Column naming confusion & 1 & 12.5\% \\
\bottomrule
\end{tabular}
\caption{Classification of semantic failures}
\label{tab:failures}
\end{table}

\section{Discussion}

\subsection{Key Findings}

\subsubsection{Accuracy-Complexity Relationship}
Semantic accuracy demonstrates a strong negative correlation with query complexity:
\begin{itemize}
    \item Simple queries (1 table): 96.0\% accuracy
    \item Multi-table queries (2-3 joins): 86.7\% accuracy
    \item Complex queries (4+ joins, subqueries): 80.0\% accuracy
\end{itemize}

This pattern aligns with observations in neural semantic parsing, where branching complexity increases error propagation.

\subsubsection{Performance Parity}
MCP-generated queries maintain comparable latency to baselines, with a mean overhead of 3.8\%. Critical for production deployment, this suggests the model learns reasonable execution strategies corresponding to actual data distributions.

\subsubsection{Optimization Heterogeneity}
Query plan metrics reveal the MCP has not learned consistent optimization principles:
\begin{itemize}
    \item 20.7\% of queries generate lower-cost plans
    \item 48.3\% of queries generate higher-cost plans
    \item Cardinality estimation accuracy is high (71.0\% exact matches)
\end{itemize}

The disparity suggests the model captures projection cardinalities but fails to optimize join orderings and predicate selectivity.

\subsection{Implications for Production Deployment}

\subsubsection{Recommendation Level}
Given 88.6\% semantic accuracy, MCP-generated queries are viable for:
\begin{itemize}
    \item Read-only reporting queries
    \item Data exploration with human review
    \item Tutorial and educational databases
\end{itemize}

Not recommended for:
\begin{itemize}
    \item Mission-critical transactional queries
    \item Unreviewed batch operations
    \item Real-time visualization generation without caching
\end{itemize}

\subsubsection{Risk Mitigation Strategies}
\begin{enumerate}
    \item \textbf{Query Validation}: Execute generated queries in sandboxed environments with row count validation
    \item \textbf{Multi-Model Consensus}: Generate $k$ candidate queries and compare results
    \item \textbf{Confidence Scoring}: Develop confidence metrics based on model uncertainty
    \item \textbf{Query Rewriting}: Apply rule-based optimizations post-generation
\end{enumerate}

\section{Limitations}

\begin{enumerate}
    \item \textbf{Dataset Size}: Evaluation on 70 queries; larger test sets would improve statistical power
    \item \textbf{Single Database System}: Results specific to Oracle Database; generalization to PostgreSQL, MySQL requires additional evaluation
    \item \textbf{Model Blackbox}: MCP server implementation not analyzed; detailed error analysis limited
    \item \textbf{Time Period}: Evaluation conducted on March 13, 2026; model updates may improve performance
    \item \textbf{Schema Complexity}: Test schema relatively modest; real-world enterprise schemas with 100+ tables not evaluated
\end{enumerate}

\section{Future Work}

\subsection{Short-term Improvements}
\begin{itemize}
    \item Expand test set to $n=500$ queries across multiple domains
    \item Implement confidence scoring using model log probabilities
    \item Develop automated query optimization using rule-based systems
    \item Add PL/SQL package generation capability
\end{itemize}

\subsection{Long-term Research Directions}
\begin{itemize}
    \item Fine-tune specialized SQL generation models on enterprise query logs
    \item Investigate query-to-query transfer learning for optimization
    \item Develop few-shot learning approaches for schema adaptation
    \item Build interactive query refinement systems with human-in-the-loop feedback
    \item Extend framework to stored procedures and complex business logic
\end{itemize}

\section{Conclusion}

This paper presents SQLclMCP, a comprehensive evaluation framework for assessing MCP-enhanced SQL generation in Oracle Database environments. Through empirical evaluation of 70 test queries across three complexity levels, we demonstrate that MCP-based SQL generation achieves \textbf{88.6\% semantic accuracy} with \textbf{3.8\% average latency overhead}. While accurate for simple queries (96.0\%), accuracy declines for complex multi-table queries (80.0\%), indicating challenges in join optimization and predicate handling.

Our EXPLAIN PLAN analysis reveals that generated queries maintain reasonable cardinality estimates (71.0\% exact match) but exhibit suboptimal join ordering. The framework provides a reusable template for evaluating code generation systems and offers practical guidance for production deployment with appropriate safeguards.

Future work should expand the evaluation to larger datasets, multiple database systems, and develop confidence scoring mechanisms for risk-aware deployment. The SQLclMCP suite serves as both a benchmarking tool and a foundation for improving neural SQL generation systems in enterprise environments.

\section*{Acknowledgments}

This research was conducted as part of the SQLclMCP initiative, a collaborative effort to evaluate and improve AI-assisted database query generation.

\begin{thebibliography}{99}

\bibitem{giordani2012semantic}
Giordani, A., \& Moschitti, A. (2012). Semantic parsing for factual question answering. \textit{In Proceedings of the Workshop on Semantic Parsing}.

\bibitem{zhong2017seq2sql}
Zhong, V., Xiong, C., \& Socher, R. (2017). Seq2SQL: Generating structured queries from natural language. \textit{arXiv preprint arXiv:1709.00103}.

\bibitem{yu2018typesql}
Yu, T., Li, Z., Zhang, Z., Zhang, R., \& Radev, D. (2018). TypeSQL: Knowledge-based type-aware neural SQL generation. \textit{arXiv preprint arXiv:1804.09769}.

\bibitem{selinger1979access}
Selinger, P. G., Astrahan, M. M., Chamberlin, D. D., Lorie, R. A., \& Price, T. G. (1979). Access path selection in a relational database management system. \textit{In Proceedings of the 1979 ACM SIGMOD International Conference on Management of Data} (pp. 23-34).

\bibitem{marcus2019learned}
Marcus, R., Papaemmanouil, O., et al. (2019). Learned cardinalities: Estimating correlated joins with deep learning. \textit{In Cidr}.

\bibitem{mcp2024}
Anthropic. (2024). Model Context Protocol: Enhancing LLM Capabilities with Contextual Information. \textit{Technical Report}.

\end{thebibliography}

\appendix

\section{Appendix: Evaluation Results Artifacts}
\label{app:artifacts}

Complete evaluation results are stored in:
\begin{verbatim}
experiments/results/mcp_evaluation_<timestamp>.json
experiments/results/mcp_evaluation_<date>_graph_<n>_*.png
\end{verbatim}

Review the generated PNG figures for detailed per-query metrics and visualizations.

\section{Appendix: Test Questions Sample}
\label{app:test_questions}

Sample test questions included in the evaluation:

\begin{lstlisting}[caption={Sample test questions from test\_questions.json}]
{
  "test_questions": [
    {
      "test_id": 1,
      "question": "Find all employees with salary greater than 50000",
      "expected_sql": "SELECT * FROM employees WHERE salary > 50000",
      "complexity": "simple"
    },
    {
      "test_id": 35,
      "question": "List employee names, departments, and salaries for sales department",
      "expected_sql": "SELECT e.name, d.dept_name, e.salary FROM employees e JOIN departments d ON e.dept_id = d.id WHERE d.dept_name = 'Sales'",
      "complexity": "medium"
    }
  ]
}
\end{lstlisting}

\end{document}
